% \section{Final Specifications and Requirements}
% This section discusses technical or methodological Requirements, data or resource requirements (software), etc.

% \section{Societal Impact}
% Study and report on project impact on society, health, safety, legal, and cultural aspects.

% \section{Environmental Impact}
% Study and report on project sustainability and environmental impact.

% \section{Ethical Issues}
% Identify ethical issues and professional responsibilities related to the design solution.

% \section{Standards - if applicable}
% Write about any standard maintained.

\section{Project Management Plan}

The project is organised into three phases over 46 weeks. The schedules are shown
in Figures~\ref{fig:plan-phase1},~\ref{fig:plan-phase2}, and~\ref{fig:plan-phase3}.

The phases follow the order of the research process. Phase 1 establishes the physical
and methodological foundation. Phase 2 develops and tests the algorithms under
controlled conditions. Phase 3 applies the validated pipeline to real observations
and completes the reporting.

This order reduces the risk of interpreting real data with an untested method.

\vspace{0.5cm}
\textbf{Phase 1 (Weeks 1 to 15). Background and Method Design}

The first 15 weeks are reserved for background study and method design. The later
experiments depend on both of them.

The literature review is completed first so that the project requirements and research
gaps are clear. The clustering and causal discovery reviews come next. They are
followed by the accretion-disk review, which connects the statistical problem to the
physical variables.

The gaps found in these reviews guide the pipeline design. Data collection or
simulation is scheduled at the end of this phase. The data generator and preprocessing
choices should follow the final model design.

\vspace{0.5cm}
\begin{enumerate}
    \item Review causal discovery methods for time series.
    \item Review unsupervised methods for regime identification.
    \item Review accretion-disk physics and known accretion states.
    \item Identify gaps in causal discovery for astrophysical time series.
    \item Design the regime-separation and graph-recovery pipeline.
    \item Obtain or simulate accretion-disk time-series data from sources such as GRMHD simulations or NICER/RXTE archives.
\end{enumerate}

\vspace{0.5cm}
\begin{figure}[htbp]
\centering
\begin{center}
\begin{ganttchart}[
    x unit=0.8cm,
    y unit title=0.8cm,
    y unit chart=0.8cm,
    vgrid,
    hgrid,
    title label font=\bfseries\small,
    bar label font=\small,
    bar/.style={fill=blue!20},
    include title in canvas=false
]{1}{15}
\gantttitle{Phase 1. Weeks 1 to 15}{15} \\
\gantttitlelist{1,...,15}{1} \\
\ganttbar{Task 1}{1}{3} \\
\ganttbar{Task 2}{4}{6} \\
\ganttbar{Task 3}{7}{9} \\
\ganttbar{Task 4}{10}{11} \\
\ganttbar{Task 5}{12}{14} \\
\ganttbar{Task 6}{14}{15}
\end{ganttchart}
\end{center}
\caption{Project plan for Phase 1.}
\label{fig:plan-phase1}
\end{figure}
\vspace{0.5cm}

Figure~\ref{fig:plan-phase1} shows that the tasks are arranged from general to
specific. The first reviews establish the research context. The physics review and
gap analysis then define the scientific and technical requirements.

Pipeline design comes next, followed by data preparation. This sequence ensures that
the data and experiments match the proposed method rather than being selected
independently of it.

\vspace{0.5cm}
\textbf{Phase 2 (Weeks 16 to 30). Algorithm Development and Simulation}

Phase 2 uses 15 weeks to implement and test the method before real-data analysis.
The first step is to generate data with known regimes and graphs. Ground truth is
required for objective evaluation.

Regime clustering is implemented next. Graph recovery within each estimated regime
comes after it. Once both components are available, the complete pipeline can be
compared with the known ground truth.

The final weeks are used for sensitivity analysis. The method is tested under
different noise levels, sample sizes, and sampling patterns.

\vspace{0.5cm}
\begin{enumerate}
    \item Generate data with known causal graphs and regime labels.
    \item Implement regime clustering based on mechanism changes.
    \item Implement time-series causal discovery within each regime.
    \item Compare the recovered regimes and graphs with the ground truth.
    \item Test sensitivity to noise, sample size, and irregular sampling.
\end{enumerate}

\vspace{0.5cm}
\begin{figure}[htbp]
\centering
\begin{center}
\begin{ganttchart}[
    x unit=0.8cm,
    y unit title=0.8cm,
    y unit chart=0.8cm,
    vgrid,
    hgrid,
    title label font=\bfseries\small,
    bar label font=\small,
    bar/.style={fill=blue!20},
    include title in canvas=false
]{16}{30}
\gantttitle{Phase 2. Weeks 16 to 30}{15} \\
\gantttitlelist{16,...,30}{1} \\
\ganttbar{Task 1}{16}{18} \\
\ganttbar{Task 2}{19}{21} \\
\ganttbar{Task 3}{22}{25} \\
\ganttbar{Task 4}{26}{28} \\
\ganttbar{Task 5}{29}{30}
\end{ganttchart}
\end{center}
\caption{Project plan for Phase 2.}
\label{fig:plan-phase2}
\end{figure}
\vspace{0.5cm}

Figure~\ref{fig:plan-phase2} reflects this dependency. Data generation comes before
algorithm testing, and regime discovery comes before regime-specific graph learning.
Validation is scheduled only after the full pipeline is available.

Sensitivity analysis is last because it evaluates the stability of the completed
implementation.

\vspace{0.5cm}
\textbf{Phase 3 (Weeks 31 to 46). Real-Data Application and Reporting}

The final 16 weeks are used for real-data application and communication of the
results. Real light curves are analysed only after the synthetic tests. The synthetic
tests show that the method can recover known structures.

The discovered regimes are then compared with known accretion states. Their graphs
are examined for physical differences. The report is drafted after the main analysis.
It then reflects the final results and limitations.

Visualisation is completed at the end. It presents the regime assignments, graphs,
and time-series results consistently throughout the thesis.

\vspace{0.5cm}
\begin{enumerate}
    \item Apply the pipeline to real accretion-disk light curves, such as those from GRS 1915+105 or Cygnus X-1.
    \item Interpret the discovered regimes using known accretion states.
    \item Compare the causal graphs and discuss their physical meaning.
    \item Write the report, including the method, results, and limitations.
    \item Prepare causal graphs, regime assignments, and time-series plots.
\end{enumerate}

\vspace{0.5cm}
\begin{figure}[htbp]
\centering
\begin{center}
\begin{ganttchart}[
    x unit=0.75cm,
    y unit title=0.8cm,
    y unit chart=0.8cm,
    vgrid,
    hgrid,
    title label font=\bfseries\small,
    bar label font=\small,
    bar/.style={fill=blue!20},
    include title in canvas=false
]{31}{46}
\gantttitle{Phase 3. Weeks 31 to 46}{16} \\
\gantttitlelist{31,...,46}{1} \\
\ganttbar{Task 1}{31}{34} \\
\ganttbar{Task 2}{35}{37} \\
\ganttbar{Task 3}{38}{40} \\
\ganttbar{Task 4}{41}{44} \\
\ganttbar{Task 5}{45}{46}
\end{ganttchart}
\end{center}
\caption{Project plan for Phase 3.}
\label{fig:plan-phase3}
\end{figure}
\vspace{0.5cm}

Figure~\ref{fig:plan-phase3} shows this final progression. Real-data analysis begins
the phase. Interpretation follows the discovery of the regimes. Reporting begins
after the main findings are available.

The final visualization task supports the written discussion. It also prepares the
project for review and submission.

\section{Risk Management}

\vspace{0.3cm}
Real astrophysical time series, including X-ray light curves, can be noisy, irregular,
and limited in size. These properties make causal graph recovery and regime
separation difficult.

Unlike in a controlled experiment, the accretion state for each observation is not known.
The number of causal mechanisms is also uncertain. Unobserved variables, such as
magnetic-field changes, may confound the measured features.

Many causal discovery algorithms also assume acyclic relationships. Accretion disks
can contain feedback among temperature, density, and accretion rate.

The study manages these risks in three ways. First, it uses simulated data with known
graphs for controlled validation. Second, it uses methods designed for time series and
time-delayed effects. Third, it compares results across methods and tests their
stability under noise and small sample sizes.

\vspace{0.5cm}
\textbf{Risks of Synthetic Data}

Synthetic data is central to Phase 2. It is the only practical way to know the
correct regime boundaries and causal graphs in advance. Without a known answer, it is
impossible to say whether the pipeline recovered the true structure. It could also
have simply produced a plausible-looking result.

Synthetic data carries its own set of risks. These risks must be stated clearly.

The first risk is that the generator may be too favourable to the method. The
simulated series are built from the same linear and piecewise-stationary assumptions
that LIRA uses to recover them. If the real system departs from these assumptions,
the simulation results will overestimate performance.

To reduce this risk, the study varies noise levels, sample sizes, and sampling
patterns. It also compares the pipeline against established baselines rather than
against an empty model.

The second risk is that synthetic data does not contain every feature of real
observations. Real X-ray light curves contain gaps caused by telescope visibility.
They also contain instrumental artefacts and source-specific variability. A simple
generator does not reproduce these effects.

A method that works perfectly on clean simulations may fail when they are present.
The sensitivity analysis therefore includes irregular and noisy conditions. It cannot
cover every case that real data may present.

The third risk is overconfidence. Because the boundary and graph are known during
simulation, it is tempting to treat a successful recovery as proof that the method
works in general.

The study avoids this. Synthetic results are treated as a necessary but insufficient
condition. The pipeline must first succeed under controlled conditions. It is only
then applied to real observations where the truth is unknown.

Any claim from the synthetic experiments is a claim about the implementation. It is
not a claim about black hole physics.

The fourth risk is that the assumed collider structure may not hold for every
system. The identifiability analysis depends on this structure. If real accretion
variables form causal chains instead of the collider used here, the identifiability
guarantee does not transfer.

This limitation is carried forward into the discussion of real data in Chapter 6. It
is not hidden by the simulation.

This risk strategy also explains the phase order. Phase 1 reduces design risk through
background study and method specification. Phase 2 reduces implementation risk by
testing the pipeline where the correct answer is known. The synthetic-data risks
above are contained and documented in this phase.

Phase 3 addresses the remaining scientific risk. It applies the validated method to
real observations. It also reports the assumptions and uncertainties that cannot be
removed from observational data.

The phase order is therefore not only a schedule. It is a deliberate escalation from
low-risk controlled tests to high-risk real-data analysis. No stage begins before the
previous stage has established its own reliability.

% \section{Economic Analysis}
% Conduct economic analysis and cost-benefit estimation.
